
\begin{proposition}[All positive canonical-flat central log spaces over a disc]\label{cdp:st:C.2.UN.71}
The positive canonical-flat central fibres of \ref{cdp:st:C.2.UN.68},
with their actual divisorial log structure and marked
dual complex, have the following complete description.
Take a rank-two lattice $N$, a finite-index sublattice
$\Gamma\subset N$, a $\Gamma$-periodic unimodular
triangulation $\mathscr P$ of $N_{\R}$ with vertex set
$N$, and a homomorphism
\[
       c:\Gamma\longrightarrow T_N
          :=\operatorname{Hom}(N^*,\C^*).
\tag{C.TG.1}\label{cdp:eq:C.TG.1}
\]
Require that each triangle have three distinct vertices
in $N/\Gamma$; this is the embedded simple-NC condition.
The central component at $p\in N/\Gamma$ is the actual
toric surface of the star fan
$\{\R_{\ge0}(\sigma-p):p\in\sigma\in\mathscr P\}$.
The homomorphism in \eqref{cdp:eq:C.TG.1} supplies all the remaining
logarithmic gluing constants, modulo the explicit torus
changes of the local charts described below.

Every such datum gives a smooth proper analytic
degeneration of connected complex two-tori over a
sufficiently small disc, with trivial canonical line,
and with exactly the specified marked central log space.
Every actual fibre in \ref{cdp:st:C.2.UN.68} is obtained in this way.
In a basis of $N$, let the columns of the integral
matrix $Q$ be a basis of $\Gamma$, and choose logarithms
of its two values under $c$ as the columns of $C$.
The actual nearby torus is
\[
 F_t=\C^2\big/
       \left(2\pi i\Z^2+(\log t\,Q+C)\Z^2\right),
       \qquad 0<|t|<\epsilon .
\tag{C.TG.2}\label{cdp:eq:C.TG.2}
\]
The central Euler number is $2[N:\Gamma]$.
This is a proper local occurrence and central-log-space
classification. It does not identify the full analytic
germ of an arbitrary given degeneration with this
chosen quotient, and it makes no compact-base
globalization assertion.
\end{proposition}

\begin{proof}
\par\medskip\noindent\textbf{The actual integral-affine charts.}
On each toric component $V_i$ of \ref{cdp:st:C.2.UN.68} write
$N_i=\pi_1(V_i\setminus D_i)$, and let $n_{ij}$
be the primitive angular meridian of its boundary
curve $C_{ij}$. A triple point $ijk$ gives the
unimodular triangle
\[
                 0,\ n_{ij},\ n_{ik}\quad\hbox{in }(N_i)_{\R}.
\tag{C.TG.3}\label{cdp:eq:C.TG.3}
\]
To pass from the chart at $j$ to that at $i$,
use the unique integral-affine map which sends
$0$ to $n_{ij}$ and whose linear part $P_{ij}$ satisfies
\[
 P_{ij}n_{ji}=-n_{ij},\qquad
 P_{ij}n_{jk}=n_{ik}-n_{ij}.
\tag{C.TG.4}\label{cdp:eq:C.TG.4}
\]
These are the actual angular maps. Indeed in a
triple chart $f=xyz$, the coordinates on $V_i$
are $(y,z)$ and those on $V_j$ are $(x,z)$.
On the nearby fibre $x=t/(yz)$, a loop varying
$x$ with $z$ fixed has angular vector $-n_{ij}$,
and a loop varying $z$ with $x$ fixed has vector
$n_{ik}-n_{ij}$. Multiplication of these coordinates
by holomorphic units leaves these integral maps
unchanged after the chart is shrunk.

There are exactly two triple points on $C_{ij}$.
Write their third vertices as $k,l$, retaining
their two distinct incidences if the same global
component occurs twice. The toric fan relation on
$V_i$ is
\[
 n_{ik}+n_{il}=-(C_{ij}^2)_{V_i}n_{ij}.
\tag{C.TG.5}\label{cdp:eq:C.TG.5}
\]
The corresponding relation at $j$ agrees with
\eqref{cdp:eq:C.TG.4} at the other endpoint precisely because
\[
             (C_{ij}^2)_{V_i}+(C_{ij}^2)_{V_j}=-2,
\tag{C.TG.6}\label{cdp:eq:C.TG.6}
\]
the actual normal-product and two-triple-point
formula of \ref{cdp:st:C.2.UN.65}. Thus there is no affine singularity
in an edge. The complete smooth toric fan at each
vertex makes its star an affine neighbourhood of
the origin. The triangle charts and \eqref{cdp:eq:C.TG.4}
therefore give an integral-affine structure on
the entire actual dual two-torus $B$.

\par\medskip\noindent\textbf{The angular lattices inject into the actual nearby torus.}
The nowhere-zero three-form of \ref{cdp:st:C.2.UN.65} is, in an
actual triple chart,
\[
       \Omega=a(x,y,z)\,dx\wedge dy\wedge dz,
          \quad a(0)\ne0,
          \qquad f=xyz.
\tag{C.TG.7}\label{cdp:eq:C.TG.7}
\]
The induced relative two-form on $F_t$ is
$a(x,y,t/(xy))\,d\log x\wedge d\log y$.
Its integral over the real angular two-torus
$|x|=|y|=|z|=|t|^{1/3}$ is exactly
\[
 (2\pi i)^2\sum_{q\ge0}a_{q,q,q}t^q
                  =(2\pi i)^2(a(0)+O(t))\ne0.
\tag{C.TG.8}\label{cdp:eq:C.TG.8}
\]
This follows by the convergent power-series expansion
and the two circle integrals; only the diagonal
monomials survive.

For a fixed component, its toric complement retracts
onto such an angular two-torus. Any finite collection
of its loops, their homotopies and the connecting
paths to the corner lies in a compact subset away
from the boundary. The submersion $f$ near this
compact subset gives actual smooth horizontal
transports into $F_t$. Consequently its angular
lattice has a well-defined homomorphism
$N_i\to\pi_1(F_t)$, up to conjugation, and the
corner homomorphism is this same map. Since
$\pi_1(F_t)=\Z^4$ is abelian, conjugation makes
no ambiguity. If this map had rank at most one,
the fundamental class of the angular torus would
map to zero in $H_2(F_t,\Z)$: the actual complex
torus is a $K(\Z^4,1)$ and its degree-two homology
map is the exterior square of its degree-one map.
This contradicts \eqref{cdp:eq:C.TG.8}. Thus every $N_i$
injects into the actual $\Z^4$.

The corner maps and the compact double-curve collars
identify the injected lattices by exactly \eqref{cdp:eq:C.TG.4}.
Transport around any loop in $B$ gives a linear
holonomy $A$ with $\iota_iA=\iota_i$ in $\Z^4$.
Injectivity gives $A=I$. This proves trivial linear
holonomy by maps into the given smooth fibre; no
invariant-cycle surjectivity is being assumed.

The affine charts now differ by translations.
Their Euclidean metrics give a complete flat metric
on compact $B$. Extending straight lines on its
universal cover proves that the developing map is
an actual affine isomorphism with $N_{\R}$.
The deck group is a full lattice of translations
$\Gamma\subset N$, because all the transition
translations in \eqref{cdp:eq:C.TG.3}--\eqref{cdp:eq:C.TG.4} are integral.
Every developed triangle is unimodular. Conversely
every lattice point of $N$ is a vertex, since a
unimodular triangle contains no other lattice point
and the triangles cover $N_{\R}$. Thus the actual
dual complex is $(N_{\R},\mathscr P)/\Gamma$,
with exactly $[N:\Gamma]$ vertices.

\par\medskip\noindent\textbf{The exact logarithmic gluing comparison.}
We spell out the primary-source comparison used to
retain the given gluing constants. The divisorial
log structure of $(X,S)$ restricted to $S$ is locally
generated by the actual factors of $f$. At a triple
point its chart is $\mathbb N^3$, and the base generator
is $(1,1,1)$. At a double curve it is the analogous
two-factor chart. The normalization is the disjoint
union of the actual toric $V_i$; its conductor is
the reduced full toric boundary, with the exact
cartesian and cocartesian SNC conductor square.
The central surface is Gorenstein and this log
map is smooth at every point, directly from
$xyz=t$, $xy=t$ and the submersion chart. Thus
the log-singular set is empty.

The precise gluing result is
\cite[Theorems 4.14, 5.2 and 5.4]{GrossSiebert2006I}.
Its use here does not import a polarization or an
algebraic total space. The following description
checks the analytic-to-algebraic comparison in the
proof. Fix the toric identifications of every
normalization stratum with its actual counterpart.
Every boundary identification is an automorphism
of $\PP^1$ preserving its two specified corner
points, and hence, in the specified coordinates,
is multiplication by a nonzero constant. These
are the closed gluing maps of that source.

At a triple point, restrict an actual log chart
to each of the three toric affine planes. A
monomial $z^p$ is multiplied by a unit $h_p$
in its analytic local ring, with
$h_{p+q}=h_ph_q$. The map
\[
       p\longmapsto h_pz^p
       \quad\longmapsto\quad
       p\longmapsto h_p(0)z^p
\tag{C.TG.9}\label{cdp:eq:C.TG.9}
\]
is exactly the units-to-residues map in the proof
of Theorem 4.14. Restricting to a shared coordinate
axis and evaluating at its endpoint commute.
The resulting constant monoid maps therefore
intertwine the actual closed gluing diagrams.
Their colimits give the actual open toric charts
and the isomorphism with the given central space,
as in the commutative squares of that proof.
All these maps are defined by the same convergent
local rings in the analytic setting. The constant
residues are algebraic monomial maps; thus the
resulting finite open atlas and its actual
equivalence relation also give an algebraic
space whose analytification is the given $S$.

The remaining log data are the nowhere-zero
sections of the lines
$N_{C/V_i}\otimes N_{C/V_j}\otimes\OO_C(p+q)$
on the double $\PP^1$'s, where $p,q$ are their
two triple points, with the actual triple
compatibilities. These corrected normal-product
lines have degree zero by \eqref{cdp:eq:C.TG.6}, and
their sections are nonzero constants in their
trivializations. They are algebraic as well.
This is the exact log-smooth-structure torsor
in \cite[Example 3.30 and Theorem 3.22]{GrossSiebert2006I};
the constant maps just described retain its
given section, not merely the underlying surface.
Consequently the log central space satisfies
Definition 4.3 of that source over $\C$ with
empty singular set. Positivity says these
sections have no poles and is immediate.
The affine monodromy polytopes are points;
in Definition 1.60 simplicity is the case
$p=0$, whose remaining conditions are vacuous.

Theorems 5.2 and 5.4 now supply actual torus
changes of the open charts, preserving the
marked dual complex, which turn these data
into a Cech cocycle for the constant sheaf
$N\otimes\C^*=T_N$ on $B$.
The maps of a cohomologous change are explicitly
$s'_{\alpha\beta}=a_\alpha^{-1}
s_{\alpha\beta}a_\beta$, acting on each chart
by its torus element $a_\alpha$; their pullbacks
also preserve the normalized log section.
Thus they are actual isomorphisms of central
log spaces. On the universal cover
$N_{\R}$ a cocycle is trivialized by path
products, and its only remaining values are
the two commuting deck holonomies. Therefore
\[
 H^1(B,N\otimes\C^*)
       =\operatorname{Hom}(\Gamma,T_N),
\tag{C.TG.10}\label{cdp:eq:C.TG.10}
\]
with precisely the marked gluing datum \eqref{cdp:eq:C.TG.1}.
This is the actual comparison to be used in
the construction, rather than an assertion
that arbitrary analytic neighbourhoods are toric.

\par\medskip\noindent\textbf{The smooth toric threefold and its exact deck maps.}
For every triangle $\sigma$ of $\mathscr P$, take
the cone over $\sigma\times\{1\}$ in
$N_{\R}\oplus\R$. These cones and their faces
form a fan $\Sigma$. Its three primitive
generators $(p_0,1),(p_1,1),(p_2,1)$ are an
integral basis because the triangle is
unimodular. Its toric chart is therefore
literally $\C^3$ with coordinates $x_0,x_1,x_2$,
and the character of the last dual coordinate is
\[
                  t=x_0x_1x_2,
       \qquad z^m=\prod_{a=0}^2x_a^{\langle m,p_a\rangle}
                       \quad(m\in N^*).
\tag{C.TG.11}\label{cdp:eq:C.TG.11}
\]
The charts glue by their actual Laurent
monomial localizations. The fan intersection
condition makes this a Hausdorff smooth toric
analytic threefold. Although the fan has
infinitely many cones, each central orbit has
the ordinary finite star of its cone, and the
open torus is a single common open chart.
These are actual analytic neighbourhoods.

A translation $\lambda\in\Gamma$ acts on the
lattice $N\oplus\Z$ by the integral shear
$(n,h)\mapsto(n+h\lambda,h)$. It preserves
the fan and induces the actual monomial map
between its charts. Multiply this map by
$c_\lambda$ on the torus factor. On the
punctured open torus the result is exactly
\[
              g_\lambda(z,t)=(c_\lambda t^\lambda z,t).
\tag{C.TG.12}\label{cdp:eq:C.TG.12}
\]
The $c_\lambda$ are a homomorphism, and
constant torus multiplication commutes with
the shears. Thus $g_\lambda g_\mu=g_{\lambda+\mu}$
on every chart. On the central fibre these
maps identify the component and boundary
charts by the cocycle in \eqref{cdp:eq:C.TG.10}. The
quotient central log space is consequently
the exact space specified earlier, including
its given log section.

\par\medskip\noindent\textbf{Freeness, properness and the actual compact fibres.}
Extend the homomorphism
$\lambda\mapsto\log|c_\lambda|$ real-linearly
to a map $K:N_{\R}\to N_{\R}$.
For $0<|t|<\epsilon$ put $L=\log|t|$ and
\[
                 q=\frac{\log|z|}{L}.
\tag{C.TG.13}\label{cdp:eq:C.TG.13}
\]
Then \eqref{cdp:eq:C.TG.12} sends $q$ to
$q+(I+K/L)\lambda$.
Choose $\epsilon<1$ so small that $I+K/L$
and its inverse are uniformly bounded and
invertible for $0<|t|<\epsilon$.
The action on every nonzero fibre is
therefore free and properly discontinuous;
its real logarithmic quotient has a compact
fundamental parallelogram. Its angular
directions already form a compact two-torus.
Taking logarithms gives exactly \eqref{cdp:eq:C.TG.2},
whose four real periods are independent
because $LQ+\operatorname{Re}C$ is invertible.

At $t=0$ a nonzero translation cannot fix
any orbit: it would preserve a bounded cell
of the triangulation, which is impossible.
The central action is thus free. To verify
proper discontinuity also at the central
fibre, use a finite collection of fixed
toric charts about a proposed compact set.
For points in those charts tending to zero,
the coordinates in \eqref{cdp:eq:C.TG.11} show that $q$
stays bounded: each $\log|x_a|/L$ is bounded
below by a quantity tending to zero and
the three coefficients sum to one.
If both a point and its $g_\lambda$-translate
stay in such fixed compact sets, \eqref{cdp:eq:C.TG.13}
and the uniform inverse bound force
$\lambda$ to stay bounded. Only finitely
many lattice elements occur. The same
argument at central points uses the finite
sets of their translated cells. Hence the
quotient is Hausdorff and smooth.

Finally choose the representatives of $q$
in a uniformly bounded fundamental
parallelogram for $(I+K/L)\Gamma$.
Only finitely many triangles of
$\mathscr P$ meet this bounded set.
In a triangle containing $q$, its
barycentric coordinates are precisely
\[
             \ell_a=\frac{\log|x_a|}{L},
             \qquad \ell_a\ge0,\quad\sum_a\ell_a=1.
\tag{C.TG.14}\label{cdp:eq:C.TG.14}
\]
Consequently $|x_a|\le1$ in that chart.
Any sequence of fibres tending to zero
has, after its actual deck translations,
a subsequence in one of these finitely
many closed unit polydiscs, and hence a
convergent subsequence. The central fibre
itself is the finite union of the complete
toric star surfaces. This proves properness
over a smaller disc. The punctured fibres
are connected by their literal quotient
presentation, and the central dual torus
is connected as well.

The triangle-distinctness condition in
the statement prevents a component from
meeting a translate of itself. The
quotient therefore retains smooth embedded
components and the full reduced
simple-NC equation in \eqref{cdp:eq:C.TG.11}.
The differential
\[
                  dt\wedge d\log z_1\wedge d\log z_2
\tag{C.TG.15}\label{cdp:eq:C.TG.15}
\]
is invariant under \eqref{cdp:eq:C.TG.12}, because all
additional terms are multiples of $dt$.
In each unimodular chart \eqref{cdp:eq:C.TG.11} it is,
up to the basis determinant $\pm1$,
$dx_0\wedge dx_1\wedge dx_2$.
It extends nowhere vanishingly across
the central fibre. Thus the constructed
total space has trivial canonical line.
The exact central-space comparison and
\ref{cdp:st:C.2.UN.68} give its Euler number and complete
both the occurrence and exhaustion
assertions in the stated local scope.
\end{proof}


